\documentclass[12pt, a0paper]{article}

% ===== 页面 =====
\usepackage[
    paperwidth=84.1cm,
    paperheight=118.9cm,
    margin=2cm,
    columnsep=2cm
]{geometry}

% ===== 中文 =====
\usepackage{ctex}

% ===== 双栏 =====
\usepackage{multicol}

% ===== 表格 =====
\usepackage[table]{xcolor}
\usepackage{colortbl}
\usepackage{multirow}
\usepackage{booktabs}
\usepackage{tabularx}

% ===== 其他 =====
\usepackage{tikz}
\usepackage{enumitem}
\usepackage{amsmath}
\usepackage{amssymb}

% ===== 颜色 =====
\definecolor{myblue}{RGB}{25,55,109}
\definecolor{myorange}{RGB}{200,100,20}
\definecolor{mygreen}{RGB}{30,120,70}
\definecolor{myred}{RGB}{180,40,40}
\definecolor{mygray}{RGB}{130,130,130}
\definecolor{mypurple}{RGB}{100,40,160}
\definecolor{lightblue}{RGB}{235,245,255}

% ===== 字体缩放 (A0) =====
\renewcommand{\normalsize}{\fontsize{22}{30}\selectfont}
\renewcommand{\large}{\fontsize{28}{36}\selectfont}
\renewcommand{\Large}{\fontsize{34}{42}\selectfont}
\renewcommand{\LARGE}{\fontsize{40}{48}\selectfont}
\renewcommand{\huge}{\fontsize{48}{56}\selectfont}
\renewcommand{\Huge}{\fontsize{56}{64}\selectfont}
\renewcommand{\small}{\fontsize{18}{24}\selectfont}
\renewcommand{\footnotesize}{\fontsize{15}{20}\selectfont}

% ===== 收紧间距 =====
\setlength{\parskip}{0.3em}
\setlength{\parindent}{0pt}
\setlength{\multicolsep}{0.5em}
\raggedbottom

% ===== 节标题 =====
\usepackage{titlesec}
\titleformat{\section}[hang]
    {\huge\bfseries\color{myblue}}
    {\thesection}
    {1em}
    {}
    [\vspace{-1mm}{\color{myblue}\rule{\columnwidth}{1.5pt}}]
\titlespacing{\section}{0pt}{0.5em}{0.3em}

\titleformat{\subsection}[hang]
    {\Large\bfseries\color{myblue!80}}
    {\thesubsection}
    {1em}
    {}
\titlespacing{\subsection}{0pt}{0.3em}{0.1em}

% ===== 列表 =====
\setlist{nosep, leftmargin=2.5em, itemsep=0.3em}

% ===== 占位图 =====
\newcommand{\placeholder}[2]{%
    \begin{center}
    \colorbox{mygray!12}{\parbox{0.9\columnwidth}{%
        \centering
        \vspace{1.5em}
        {\Large\bfseries\color{mygray} [ 占位图： #1 ]}
        \vspace{0.8em}
        {\small\color{mygray} #2}
        \vspace{1.5em}
    }}
    \end{center}
}

\begin{document}

% =====================
% 标题
% =====================
\begin{center}
    {\fontsize{54}{64}\selectfont\bfseries\color{myblue}
        从任务专用脑电解码器到脑电基座大模型：\\[3mm]
        EEGNet、EEG Conformer 与 CBraMod 的复现与扩展评测\\[10mm]
    }

    {\fontsize{28}{34}\selectfont
        第一作者 (SUSTech ID) \hspace{2em} 第二作者 (ID) \hspace{2em} 第三作者 (ID)\\[4mm]
    }

    {\fontsize{20}{28}\selectfont\color{mygray}
        南方科技大学生物医学工程系 \hspace{3em}
        BMEB215 机器学习与医学工程应用 · 2026 Spring
    }

    \vspace{6mm}
    {\color{myblue}\hrule height 2.5pt}
    \vspace{6mm}
\end{center}

\begin{multicols}{2}

% ================================================================
% 1. 背景 (Background)
% ================================================================
\section{研究背景 (Background)}

\subsection{为什么需要脑电基座大模型？}

EEG 解码长期面临三个核心困难：

\begin{itemize}
    \item \textbf{标注数据少} —— 多数 BCI、临床、情绪、睡眠数据集被试数量有限，标签成本高。
    \item \textbf{跨域差异大} —— EEG 信号受被试、session、设备、导联、采样率和预处理流程影响明显。
    \item \textbf{评测体系碎片化} —— 不同论文使用不同 split、滤波、输入尺度，模型之间难以公平比较。
\end{itemize}

传统任务专用模型 (CNN/RNN/Transformer) 在特定数据集上从零训练，优势是结构明确、参数量小、单任务性能强；劣势是缺少跨任务复用能力，换一个任务往往要重新设计、重新调参。

基座大模型改变的是训练范式：先通过大规模 EEG 数据预训练学习通用表征，再通过 Linear Probing (LP) 或 Fine-Tuning (FT) 适配不同下游任务。

\subsection{发展脉络}

\begin{center}
\begin{tabular}{llll}
    \toprule
    \textbf{阶段} & \textbf{代表工作} & \textbf{为什么选它} & \textbf{主要价值} \\
    \midrule
    深度学习进入 EEG & EEGNet (2018) & 紧凑、领域感知的深度监督模型 & 高效学习 EEG 特征 \\
    全局时序建模 & EEG Conformer (2023) & CNN 前端 + Transformer 注意力 & 长程依赖 + Attention 可视化 \\
    大规模预训练 & CBraMod (2025) & EEG Foundation Model & 预训练 → LP/FT 多任务适配 \\
    \bottomrule
\end{tabular}
\end{center}

\vspace{1em}

\textbf{核心问题：}EEGNet 和 EEG Conformer 代表任务专用深度模型，CBraMod 代表从"训练专家模型"到"预训练通用表征再适配下游任务"的范式转变。为什么在专家模型已能在单任务上取得强性能时，我们仍需要脑电基座大模型？

% ================================================================
% 2. 方法 (Methods)
% ================================================================
\section{模型方法 (Methods)}

\subsection{EEGNet —— 深度学习 EEG 解码代表}

用紧凑 CNN 将 EEG 领域知识融入深度学习，是 EEG 解码中最经典的专家模型。

\textbf{架构：}Temporal Convolution (频率时间滤波器) $\rightarrow$ Depthwise Spatial Convolution (通道空间模式) $\rightarrow$ Separable Convolution (降参)。每个下游任务单独监督训练。

\textbf{定位：}任务专用深度模型 baseline，与 EEG Conformer、CBraMod 形成"专家模型 vs 基座模型"对照。

\vspace{1em}

\subsection{EEG Conformer —— Transformer 增强专家模型}

在卷积局部特征提取之后引入 Transformer，用 Self-Attention 建模 Patch 间全局依赖。

\textbf{架构：}CNN 前端提取局部 EEG 特征并形成 Patch Embedding $\rightarrow$ Transformer Encoder 建模长程时序依赖 $\rightarrow$ 分类头。Attention Map 可用于可视化分析。

\textbf{定位：}比 EEGNet 更复杂的任务专用 baseline。进行"去掉 Transformer block"消融，检验 Attention 对单任务 EEG 解码的贡献。

\vspace{1em}

\subsection{CBraMod —— EEG 基座大模型}

通过大规模预训练获得通用 EEG 表征，再迁移到不同下游任务。

\textbf{下游适配：}CBraMod Backbone 输出 Channel-Patch 表征 $B \times C \times S \times D$。LP 冻结预训练 Backbone，只训练线性分类头；FT 同时更新 Backbone 和分类头。LP 衡量预训练表征本身是否可读，FT 衡量最终下游适配能力。

\vspace{1em}

\subsection{NeuroGate —— 神经先验融合}

CBraMod 的 Embedding 虽然泛化能力强，但高维隐空间不直接暴露可解释生理量。NeuroGate 用显式 EEG 知识补充 Foundation Embedding。

\textbf{做法：}提取人工神经先验特征 $m$（时域统计、Band Power、频谱形状、复杂度等），将 CBraMod 输出的 Backbone Embedding 记为 $h$，通过 Gated Fusion 学习门控权重 $g$：

\[
z = h + g \cdot \phi(m)
\]

目的不是让人工特征替代 Foundation Embedding，而是让模型在需要时利用可解释的 EEG 生理知识。

\placeholder{模型架构对比}{%
    EEGNet / EEG Conformer / CBraMod / NeuroGate 结构示意图。\\
    展示四类模型的关键模块和架构差异。\quad (等待绘制)}

% ================================================================
% 3. 结果 (Results)
% ================================================================
\section{实验结果 (Results)}

\subsection{CBraMod 下游任务能否复现？}

{\footnotesize
\begin{center}
\rowcolors{2}{lightblue}{white}
\begin{tabularx}{\columnwidth}{l>{\raggedright\arraybackslash}p{3.8cm}ccccccc}
\toprule
\multirow{2}{*}{\textbf{数据集}} & \multirow{2}{*}{\textbf{版本配置}} & \multirow{2}{*}{\textbf{Seeds}} &
\multicolumn{2}{c}{\textbf{bAcc}} & \multicolumn{2}{c}{\textbf{\(\kappa\)}} &
\multirow{2}{*}{\textbf{论文参考}} & \multirow{2}{*}{\textbf{结论}} \\
\cmidrule(lr){4-5} \cmidrule(lr){6-7}
& & & Mean & Best & Mean & Best & & \\
\midrule
BCIC2A & session T/E+CAR & 42/2024/3407 & 0.6382 & 0.6709 & 0.5177 & 0.5612 & \(\kappa\) 0.5138 & \cellcolor{mygreen!15}对齐 \\
FACED & baseline official FT & 42/2024/3407 & 0.5245 & 0.5340 & 0.4617 & 0.4718 & bAcc 0.5509 & 接近 \\
ISRUC-S1 & baseline official FT & 42/2024/3407 & 0.7780 & 0.7850 & 0.7128 & 0.7373 & bAcc 0.7865 & \cellcolor{mygreen!15}基本复现 \\
Physionet-MI & flatten FT & 42/2024/3047 & 0.5975 & 0.6181 & 0.4631 & 0.4903 & bAcc 0.6417 & 接近 \\
SEED-V & baseline official FT & 42/2024/3407 & 0.3927 & 0.4018 & 0.2450 & 0.2592 & bAcc 0.4091 & 接近 \\
SHU-MI & flatten FT & 42/2024/3047 & 0.5851 & 0.6380 & 0.1700 & 0.2756 & bAcc 0.6370 & \cellcolor{mygreen!15}best 达到 \\
MDD & flatten FT & 42/2024/3047 & \textbf{0.9034} & \textbf{0.9674} & \textbf{0.8008} & \textbf{0.9099} & bAcc 0.9560 & \cellcolor{mygreen!15}best 超过 \\
TUEV & baseline official FT & 42/2024/3407 & 0.5499 & 0.5755 & 0.5964 & 0.6343 & \(\kappa\) 0.6671 & 未达 \\
\bottomrule
\end{tabularx}
\end{center}
\vspace{0.3em}
{\small\color{mygray}注：bAcc = 平衡准确率，\(\kappa\) = Cohen's Kappa。Mean 为三种子均值，Best 为最优种子。论文参考值口径与对应列一致。}
}

\subsection{EEGNet 复现结果}

作为任务专用专家模型 Baseline，我们使用原始 EEGNet 在四个运动想象数据集上进行了复现：

{\footnotesize
\begin{center}
\rowcolors{2}{lightblue}{white}
\begin{tabular}{lcccr}
\toprule
\textbf{数据集} & \textbf{类别数} & \textbf{Val Acc (\%)} & \textbf{Test Acc (\%)} & \textbf{Test \(\kappa\)} \\
\midrule
BCICIV2A & 4 & 79.17 & -- & -- \\
BCICIV2B & 2 & 77.72 & 78.35 & 0.5669 \\
IIIa & 4 & 55.95 & -- & -- \\
Physionet MI (22ch) & 4 & 49.30 & -- & -- \\
\bottomrule
\end{tabular}
\end{center}
\vspace{0.3em}
{\small\color{mygray}注：BCICIV2A 和 BCICIV2B 为主要结果 (300 epochs)；IIIa 和 Physionet MI 为快速基线 (100 epochs, single seed)，仅作为实现通用性的补充证据。}
}

\subsection{复现中的关键经验}

CBraMod 输出 $B \times C \times S \times D$，其中 $C$ 是通道、$S$ 是 Patch/Segment、$D$ 是隐层维度。如果下游 Head 过早对通道求平均 ($B \times C \times S \times D \rightarrow B \times S \times D$)，会丢失 MI 等任务依赖的空间分布信息。

\textbf{可讲结论：}预训练表征不等于下游性能；Foundation Model 还需要与 EEG 结构匹配的 Downstream Readout。

\placeholder{Readout 修正前后对比}{%
    SHU-MI / Physionet-MI / MDD 数据集上的修正前 vs 修正后 bAcc。\\
    保留 Channel-Patch 结构后部分任务明显恢复。\quad (等待队友数据)}

\subsection{NeuroGate：神经先验能否增强预训练表征？}

\textbf{可讲结论：}NeuroGate 的主证据在 LP，说明显式神经先验能补充 Frozen Foundation Embedding；FT 证据可作为辅助支持。详见表~\ref{tab:neurogate-lp} 和表~\ref{tab:neurogate-ft}。

{\footnotesize
\begin{center}
\rowcolors{2}{lightblue}{white}
\begin{tabular}{lccc}
\toprule
\textbf{数据集} & \textbf{CBraMod LP bAcc} & \textbf{NeuroGate LP bAcc} & \(\mathbf{\Delta}\) \\
\midrule
HMC & 0.6149 \(\pm\) 0.0256 & 0.6922 \(\pm\) 0.0172 & \cellcolor{mygreen!15}\(+\)0.0773 \\
ISRUC-S1 & 0.5962 \(\pm\) 0.0579 & 0.6934 \(\pm\) 0.0156 & \cellcolor{mygreen!15}\(+\)0.0972 \\
MDD & 0.8907 \(\pm\) 0.0477 & 0.9062 \(\pm\) 0.1101 & \(+\)0.0154 \\
Physionet-MI & 0.2719 \(\pm\) 0.0176 & 0.2875 \(\pm\) 0.0083 & \(+\)0.0156 \\
\bottomrule
\end{tabular}
\\[4pt]
{\small\color{mygray}表 1：冻结骨干网络下的线性探测结果 (LP)。NeuroGate 在睡眠任务 (HMC, ISRUC-S1) 上提升最为显著。}
\label{tab:neurogate-lp}
\end{center}
}

{\footnotesize
\begin{center}
\rowcolors{2}{lightblue}{white}
\begin{tabular}{lc}
\toprule
\textbf{实验设置} & \textbf{结果} \\
\midrule
LP channel-gated fusion 平均提升 (vs.\ backbone-only) & \cellcolor{mygreen!15}\(+\)0.0389 BA \\
LP gated fusion 正提升单元比例 (dataset-backbone) & 40/48 \\
CBraMod flatten-FT gated fusion 平均提升 & \(+\)0.0078 BA (9/12 为正) \\
FT best-of-gated/logit 平均提升 & \(+\)0.0097 BA \\
\bottomrule
\end{tabular}
\\[4pt]
{\small\color{mygray}表 2：Selected-12 数据集上 FT 实验统计结论。}
\label{tab:neurogate-ft}
\end{center}
}

\subsection{为什么需要统一 EEG Foundation Benchmark？}

目前 EEG Foundation Model 的结果分散在多篇论文中，不同模型、数据集和任务的协议不统一。统一 Benchmark 的必要性：

\begin{itemize}
    \item \textbf{公平比较} —— 同一数据集、同一 Split、同一预处理、同一指标
    \item \textbf{发现泛化边界} —— 跨任务 Benchmark 暴露哪些任务仍然困难
    \item \textbf{区分模型能力与协议差异} —— 减少因数据处理不同导致的伪差距
\end{itemize}

\placeholder{统一 Benchmark 扩展结果}{%
    HMC / MDD / Siena\_EEG / SleepEDF\_full 等数据集上的 CBraMod 表现。\\
    3-seed avg bAcc + 外部参考值对比。\quad (等待队友数据)}

% ================================================================
% 4. 总结 (Summary)
% ================================================================
\section{总结 (Summary)}

\subsection{主要结论}

\begin{enumerate}[leftmargin=2em]
    \item \textbf{EEGNet 和 EEG Conformer 是任务专用专家模型代表}，展示了从紧凑深度模型到 Transformer 增强模型的路线。
    \item \textbf{CBraMod 代表另一种范式：大规模预训练 + 下游适配}，优势不一定是单任务最高分，而是跨任务复用、少标签迁移和 Benchmark 扩展。
    \item \textbf{专家模型和基座模型应该互补评价}，专家模型衡量单任务上限，基座模型衡量通用表征和迁移能力。
    \item \textbf{CBraMod 的下游性能依赖合理 Readout}，保留 Channel-Patch 结构比过早通道池化更适合 EEG 任务。
    \item \textbf{NeuroGate 提供了可解释性补充}，通过 Gated Fusion 注入显式 EEG 神经先验，LP 证据最强。
    \item \textbf{统一 Benchmark 是 EEG Foundation Model 发展的必要基础}。
\end{enumerate}

\subsection{如果评委只记住一句话}

\begin{center}
\colorbox{lightblue}{\parbox{0.95\columnwidth}{\centering
    {\LARGE\bfseries
        EEG Foundation Models 正在解决泛化问题，\\[3mm]
        而可解释性正在成为新的瓶颈。}\\[4mm]
    {\large
        NeuroGate 尝试将神经科学知识重新引入 Foundation Models\\
        —— 其中一个可能的方向。}
}}
\end{center}

% ================================================================
% 参考文献
% ================================================================
\vspace{1em}
\section*{参考文献}

{\footnotesize
\begin{enumerate}[leftmargin=2em]
    \item Lawhern et al. \textbf{EEGNet: A Compact Convolutional Neural Network for EEG-based Brain-Computer Interfaces}. \textit{J. Neural Eng.}, 2018.
    \item Song et al. \textbf{EEG Conformer: Convolutional Transformer for EEG Decoding and Visualization}. \textit{IEEE TNSRE}, 2023.
    \item Wang et al. \textbf{CBraMod: A Criss-Cross Brain Foundation Model for EEG Decoding}. \textit{ICLR}, 2025.
    \item NeuroGate / NeuroKE 内部参考：neural-prior gated fusion experiments under LP and FT.
    \item 本地复现记录：\texttt{cbramod\_reproduction\_results\_summary.md}
\end{enumerate}
}

\end{multicols}

% =====================
% 底部
% =====================
\vspace{4mm}
{\color{myblue}\hrule height 2pt}
\vspace{3mm}
\begin{center}
{\footnotesize\color{mygray}
    BMEB215 · 机器学习与医学工程应用 · 2026 Spring \hspace{3em}
    南方科技大学 · 生物医学工程系
}
\end{center}

\end{document}
